% =====================================================================
%  frontmatter/howtoread.tex
% =====================================================================
\chapter*{How to Read This Book}
\addcontentsline{toc}{chapter}{How to Read This Book}
\markboth{How to Read This Book}{How to Read This Book}

% TODO: keep this short.  Suggested skeleton:

\section*{The tag system}
% Sections are tagged in the margin with one of the following labels.
%   Core            — belongs in a standard Calculus I or II course.
%   Bridge          — prepares ideas needed later (multivariable calculus,
%                     differential forms, differential equations, analysis).
%   Proof           — gives a more precise theorem or proof.
%   Counterexample  — shows why a theorem needs its hypotheses.
%   Computing       — uses numerical or computer-based exploration.
%   Optional        — enrichment for honors courses or projects.
%   Project         — a longer exploration, often applied or computational.

\section*{Suggested course paths}
% TODO: paste the four paths from the TOC introduction
%   - Standard Calculus I path
%   - Standard Calculus II path
%   - Honors / proof-enriched path
%   - Applied science path
